\section{Conclusion}
\label{sec:conclusion}

Sirus demonstrates that a selective x86-64 code virtualizer can be built as a
zero-dependency, safe-Rust artifact with a real evidence trail. The system lifts
a supported function into a bounded IR, flattens and rewrites its control flow,
encodes it into a per-build randomized encrypted ISA, synthesizes a freestanding
\code{no_std} VM, injects it into an ELF or PE image, and removes the original
body. The default runtime forbids unsafe code and host memory access; the
optional hardened tier isolates the controls that require an unsafe boundary.

The evaluation is deliberately mixed. The positive results are substantial: 226
tests in debug and release modes, Kani and Miri coverage of bounded components,
tens of millions of fuzzing executions, differential equality with native code
across the supported feature fixtures, external decompilers that reduce the
patched functions to stubs, bounded symbolic probes that do not find a preimage,
three external MBA simplifiers that produce a measured cost asymmetry, a
constant-time subset below the conventional timing threshold, and green Wine,
native Windows, and MSVC \code{/guard:cf} jobs. The negative results are equally
important: the default branchy profile leaks timing, the interpreter skeleton is
shared between builds, the MBA layer is not irreducible, one Ghidra decompilation
does not complete, the symbolic probes time out rather than prove
unsatisfiability, and the safe tier does not provide kernel anti-debug or a
formal Spectre proof.

Sirus is therefore best understood as an engineering prototype for asymmetric
reverse-engineering cost, not as a universal obfuscator. Its contribution is the
combination of a complete selective pipeline, a randomized encrypted virtual ISA,
a fixed-pass oblivious mode, a strict safe-Rust implementation, and a
reproducible evaluation that reports its own boundaries. The remaining work is
clear: broader ISA and platform coverage, external-call and exception support,
stronger integrity and anti-debug boundaries, formal models of the virtual ISA,
and an independent audit. Those are substantial research and engineering tasks,
and the current paper deliberately does not claim them.
